
Jilid kedua \myquote{Basic Analysis} ini dimaksudkan sebagai kelanjutan yang
mulus. Penomoran bab dilanjutkan dari bagian terakhir jilid pertama. Buku ini
bermula dari catatan saya untuk mata kuliah analisis sarjana semester kedua di
University of Wisconsin---Madison pada tahun 2012, yang kurang lebih saya ajar
dengan buku Rudin. Sebagian materi dan pembuktiannya mirip dengan Rudin,
meskipun saya berusaha memberikan lebih banyak rincian dan konteks.
Pada tahun 2016, saya mengajar mata kuliah analisis sarjana semester kedua di
Oklahoma State University, sambil mengubah dan merapikan catatan tersebut dan
kali ini menggunakannya sebagai teks utama. Sejak itu, saya telah beberapa kali
lagi mengajar mata kuliah ini, menambahkan Bab \ref{approx:chapter} (yang semula
ditulis untuk mata kuliah di Wisconsin), serta membuat banyak penyempurnaan
kecil lainnya.

Pada akhirnya saya berencana menambahkan beberapa topik lagi.
Seperti biasa, saya akan berusaha mempertahankan penomoran dalam edisi-edisi
berikutnya. Topik baru yang direncanakan akan ditambahkan sebagai bab di akhir
buku atau sebagai bagian di akhir bab yang sudah ada, dan saya akan berusaha
sekuat mungkin agar nomor latihan tidak berubah.

Sebagian besar jilid kedua ini bergantung pada bagian-bagian tidak opsional dari
Jilid I\@, meskipun beberapa bagian opsional juga digunakan.
Turunan tingkat tinggi (tetapi bukan teorema Taylor itu sendiri) digunakan
dalam \ref{sec:mvhighordders}, \ref{sec:pathind}, dan
\ref{sec:mvgreenstheorem}. Fungsi eksponensial, logaritma, dan integral tak wajar
digunakan dalam beberapa contoh dan latihan, serta digunakan secara intensif
dalam \chapterref{approx:chapter}.
%It is entirely possible to avoid using all
%the optional parts from volume I\@.

%This book is not necessarily the entire second-semester course,
%though it should have enough material for an entire semester if need be.
Rencana alternatif untuk mata kuliah dua semester adalah membahas beberapa
bagian dari jilid pertama, seperti ruang metrik, pada semester kedua, sedangkan
beberapa topik opsional dari Jilid I dibahas pada semester pertama. Menyisakan
ruang metrik untuk semester kedua menjadikan semester kedua sebagai bagian
\myquote{multivariabel} dari mata kuliah tersebut.

Beberapa jalur yang mungkin ditempuh setelah ruang metrik, bergantung pada
waktu yang tersedia, adalah:
\begin{enumerate}[1)]
\item
%\volIref{Chapter \ref*{vI-ms:chapter} from volume I}{Chapter \ref{ms:chapter}},
\ref{sec:vectorspaces}--\ref{sec:svinvfuncthm},
\ref{sec:rirect}--\ref{sec:jordansets}, \ref{sec:mvchangeofvars}
(kalkulus multivariabel, dengan fokus pada integral multivariabel).
\item
%\volIref{Chapter \ref*{vI-ms:chapter} from volume I}{Chapter \ref{ms:chapter}},
Bab \ref{pd:chapter}, Bab \ref{path:chapter},
%\ref{sec:vectorspaces}--\ref{sec:mvhighordders},
%\ref{sec:diffunderint}--\ref{sec:pathind},
\ref{sec:rirect} dan \ref{sec:iteratedints}
(kalkulus multivariabel, dengan fokus pada integral lintasan).
\item Bab \ref{pd:chapter}, Bab \ref{path:chapter}, dan Bab
\ref{mi:chapter}
(kalkulus multivariabel, integral lintasan, integral multivariabel).
\item
%\volIref{Chapter \ref*{vI-ms:chapter} from volume I}{Chapter \ref{ms:chapter}},
Bab \ref{pd:chapter}, (mungkin Bab \ref{path:chapter}), dan Bab \ref{approx:chapter}
(kalkulus diferensial multivariabel, beberapa materi analisis lanjut).
\item
Bab \ref{pd:chapter}, Bab \ref{path:chapter},
\ref{sec:complexnums},
\ref{sec:arzelaascoli},
\ref{sec:stoneweier}
(variasi yang lebih sederhana dari jalur di atas).
\end{enumerate}

%When I ran the course at OSU\@, I covered the first book minus metric spaces
%and a couple of optional sections in the first semester.
%Then, in the second semester, I covered
%most of what I skipped from volume I\@, including metric spaces, and
%took option 2) above.
